\documentclass[a4paper,10pt]{report}\input{../0.Préambule/Préambule_cours_prof}\begin{document}

\newpage
\section*{Parallélogrammes particuliers}
\addcontentsline{toc}{section}{Parallélogrammes particuliers}

\begin{exer1}
Code les longueurs égales et les angles droits, sachant que le quadrilatère est :

\begin{enumerate}[font=\color{exer1}]
\begin{tabularx}{\linewidth}{*{4}X}
\item un rectangle &
\item un losange &
\item un carré &
\item Un parallélogramme\\
\end{tabularx}
\end{enumerate}

\includegraphics[scale=0.95]{../40.Image/15. Exo-Les parallélogrammes 7.png}
\end{exer1}

\noindent \begin{minipage}[t]{10cm}
\begin{exer1}
Sans justifier, complète les étiquettes sachant que $EFGH$ est un losange et $KLMN$ est un carré tel que $KM = 7$cm.

\includegraphics[scale=0.95]{../40.Image/15. Exo-Les parallélogrammes 8.png}
\end{exer1}
\end{minipage}
\hspace{2mm}
\begin{minipage}[t]{6.5cm}
\begin{exer1}
On considère un carré $KLMN$ de centre $S$ et tel que $SM=2,7$cm

\begin{enumerate}[font=\color{exer1}]
\item Fais une figure à main levée ?
\item Quelle est la mesure de l'angle $\widehat{NSM}$ ? Pourquoi ?
\item Quelle est la longueur $NS$ ? Pourquoi ?
\end{enumerate}
\end{exer1}
\end{minipage}

\noindent \begin{minipage}[t]{8cm}
\begin{exer1}
On considère le rectangle $ABCD$.

\begin{center}
\includegraphics[scale=1]{../40.Image/15. Exo-Les parallélogrammes 9.png}
\end{center}

\begin{enumerate}[font=\color{exer1}]
\item Quelle est la longueur $AC$ ? Pourquoi ?
\item Quelle est la longueur $BD$ ? Pourquoi ?
\end{enumerate}
\end{exer1}
\end{minipage}
\hspace{2mm}
\begin{minipage}[t]{8.5cm}
\begin{exer1}
On considère le losange $EFGH$.

\begin{center}
\includegraphics[scale=1]{../40.Image/15. Exo-Les parallélogrammes 10.png}
\end{center}

\begin{enumerate}[font=\color{exer1}]
\item Quelle est la mesure de l'angle $\widehat{EFG}$ ? Pourquoi ?
\item Justifie que les droites $(HF)$ et $(EG)$ sont perpendiculaire.
\end{enumerate}
\end{exer1}
\end{minipage}




\end{document}